En un proyecto de diseño en el colegio, un grupo de estudiantes debe instalar una fuente de agua en un terreno triangular, de tal forma que la distancia desde la fuente hasta cada uno de los vértices del terreno sea exactamente la misma. Un estudiante propone ubicar la fuente en el punto donde se intersectan las medianas del triángulo, argumentando que este punto es el centro del triángulo y, por lo tanto, cumple con la condición de equidistancia.

¿La afirmación del estudiante es verdadera o falsa?

Justifique su respuesta proponiendo una solución adecuada al problema y utilizando argumentos geométricos claros. No se aceptan respuestas sin justificación.